\noindent\textbf{Dificultad:} Muy difícil (final)

\subsection{Descripción del nivel}
La contraseña la valida una máquina virtual interna. Como indica el propio reto, hay que entender el \emph{bytecode} en lugar de adivinar: identificar el intérprete, su conjunto de instrucciones y el programa que verifica la contraseña.

\subsection{Análisis estático}
% Cómo se ubicó la rutina de validación: cadenas, referencias cruzadas, funciones relevantes.
\textit{[A completar.]}

\subsection{Análisis dinámico}
% Breakpoints usados, registros/memoria observados en la comparación.
\textit{[A completar.]}

\subsection{Algoritmo de validación}
% Pseudocódigo del algoritmo reconstruido.
\textit{[A completar.]}

\subsection{Solución (keygen)}
\textit{[A completar.]}

\begin{lstlisting}[language=Python]
# Keygen - Nivel 4
# TODO: pegar el código del keygen
\end{lstlisting}

\subsection{Resultado}
Nivel en análisis. \textit{[A completar.]}
